\subsubsection{Optimal Execution Problem}

\begin{principle} \hlt{Impact-Risk Trade-Off}\\
Optimal execution balances two forces: market impact rewards slower trading, while price risk rewards faster completion. A schedule is a compromise between footprint and exposure.
\end{principle}

\begin{model} \hlt{Discrete Execution Schedule}\\
Plan a buy order of size $v^\star$ over $N$ slices. During slice $n$, trade $v_n$ shares, expected market volume is $V_n$, volatility is $\sigma_n$, and the fair price evolves as $S_{n+1}=S_n+\sigma_n\xi_{n+1}$, $\sum\limits_{n=1}^{N}v_n=v^\star$.\\
Here $\xi$ is a standard Gaussian shock. The schedule controls timing; the market supplies uncertainty.
\end{model}

\begin{model} \hlt{Temporary Impact Cost}\\
A simple execution price for slice $n$ is $\widetilde{S}_n(v_n)=S_n+\kappa\sigma_n\left(\frac{v_n}{V_n}\right)^\gamma$, $W=\sum\limits_{n=1}^{N}v_n\widetilde{S}_n(v_n)$.\\
Here $\kappa$ and $\gamma$ are stock-dependent impact parameters.\\
Participation $v_n/V_n$ is the key lever: trading more of the available volume costs more.
\end{model}

\begin{definition} \hlt{Expectation-Minimising Schedule}\\
Ignoring risk, the minimiser allocates more to high-volume and low-volatility slices: $v_n^{\mathrm{exp}} =\frac{V_n/\sigma_n^{1/\gamma}}{\sum\limits_{\ell=1}^{N}V_\ell/\sigma_\ell^{1/\gamma}}v^\star$.\\
If $V_n$ and $\sigma_n$ are constant, the schedule is uniform: $v_n=v^\star/N$.
\end{definition}

\begin{mechanism} \hlt{Risk-Only Schedule}\\
If impact is ignored and only price-risk variance matters, the optimal action is to trade immediately. Waiting leaves remaining inventory exposed to random price moves.
\end{mechanism}

\begin{model} \hlt{Mean-Variance Execution Objective}\\
Let $x_n$ be the remaining quantity after slice $n$, with $x_0=v^\star$ and $x_N=0$.\\
A compact objective is $\mathcal{C} =v^\star S_0 +\kappa\sum\limits_{n=1}^{N}\sigma_n\frac{v_n^{\gamma+1}}{V_n^\gamma} +\lambda\sum\limits_{n=1}^{N}x_n^2\sigma_n^2$.\\
The middle term is impact cost; the last term is inventory risk. Higher risk aversion $\lambda$ front-loads the schedule, while higher impact sensitivity back-loads it.
\end{model}

\begin{taxonomy} \hlt{Execution Schedule Segmentation}
\begin{enumerate}[label=\roman*.]
\setlength{\itemsep}{0pt}
\item \hlt{risk-neutral}: follows volume and volatility to reduce expected impact;
\item \hlt{risk-only}: finishes immediately to remove exposure;
\item \hlt{mean-variance}: trades along a convex remaining-quantity path;
\item \hlt{tactical overlay}: lets dark pools, SORs, alpha, and limits modify child orders inside the schedule.
\end{enumerate}
\end{taxonomy}
